%%%%%%%%%%%%%%%%%%%%%%%%%%%%%%%%%%%%%%%%%%%%%%%%%%%%%%%%%%%%%%%%%%%%%%%%%%%
%%%             DESCRICOES TEORICAS E FERRAMENTAS BASICAS               %%%
%%%%%%%%%%%%%%%%%%%%%%%%%%%%%%%%%%%%%%%%%%%%%%%%%%%%%%%%%%%%%%%%%%%%%%%%%%%

\setlength{\parskip}{0.3cm}

\chapter{Fundamentação Teórica}
\label{ch:fundamentacao}

O estudo sobre o Rabbit Code demonstra como a integração entre o pensamento
computacional (PC) e sistemas gamificados amplia o engajamento no ensino de
matrizes para o nível médio \cite{rolim2023rabbit}. Ao utilizar a gamificação
como âncora de conhecimento, o software permite a exploração interativa de
conceitos matriciais, superando barreiras de aprendizagem por meio da abstração
lúdica de problemas matemáticos \cite{rolim2023rabbit}. A progressão no jogo
exige a aplicação obrigatória dos quatro pilares do PC: a decomposição ocorre ao
fragmentar o trajeto em blocos de ação; a abstração manifesta-se na análise de
linhas e colunas para definir a melhor rota; o reconhecimento de padrões é
exercitado ao identificar movimentos para evitar armadilhas; e a criação de
algoritmos consolida-se na programação sequencial necessária para atingir os
objetivos da fase \cite{rolim2023rabbit}.

A dissertação de Silva \cite{silva2020}, intitulada Pensamento computacional:
uma análise dos documentos oficiais e das questões de Matemática dos
vestibulares. Consolida-se como estudo correlato para estabelecer o referencial
teórico ao investigar diretrizes educacionais, como a BNCC, buscando contribuir
com a discussão sobre o papel ou não do Pensamento Computacional na Educação
Básica\cite{silva2020}.
  
A autora defende que a sua inserção é fundamental visto que muitas habilidades
estão próximas de conceitos matemáticos e o processo de ensino e aprendizagem
através do Pensamento Computacional pode contribuir com o desempenho de
estudantes\cite{silva2020}.

Além da fundamentação curricular, a obra destaca-se pelo mapeamento prático
dessas competências em questões reais do ENEM e vestibulares paulistas,
avaliando potenciais de Abstração e Generalização de Padrões, Algoritmos e
Processos, Representação e Automação dos Dados, Decomposição de Problemas e
Simulação e Controle de Erros \cite{silva2020}.